\section{The quotient with the prescribed positive transport}

Keep the lattice, the oriented basis $(u_0,u_1,u_2,e_4)$, and the
cell groups $D_k$ already defined. Now take $R=P^{-1}$ as the deck
action on $X$, with the same translation $a=e_4/3$. Define
\begin{equation}\label{eq:positive-quotient}
 N^+=(X\times S^1_t\times D)/
 \bigl((x,t,s)\sim(Rx,t+1/3,\zeta s)\bigr),
 \qquad q[x,t,s]=s^3.
\end{equation}
This action is free by its fourth coordinate, just as for $N$.
Here $t$ is again the fibre circle coordinate; $q$ denotes the
coordinate on the quotient disk. Let $G_P$ be the mapping torus
with relation $[x,\tau+1]=[Px,\tau]$.

\begin{theorem}\label{thm:positive-geometry}
There is an orientation-preserving diffeomorphism
\begin{equation}\label{eq:positive-product}
 \Phi^+:N^+\longrightarrow G_P\times D,
 \qquad [x,t,s]\longmapsto([x,3t],e^{-2\pi it}s).
\end{equation}
In these coordinates $q=e^{2\pi i\tau}z^3$.
The positively oriented boundary, with fibre marking at $q=1$,
is identified with $B_M$, where $M=P\times I$, by
\begin{equation}\label{eq:positive-native-boundary}
 j:B_M\longrightarrow\partial N^+,
 \qquad [x,t,v]\longmapsto[x,t+v/3,e^{2\pi iv/3}].
\end{equation}
The fibre map and the boundary map followed by retraction are
\begin{equation}\label{eq:positive-geometric-maps}
 f^+(x,t)=[x,3t],\qquad
 k^+([x,t,v])=[x,3t+v].
\end{equation}
Thus the positive transport in this marking is exactly $T$.
\end{theorem}
\begin{proof}
Under the quotient generator,
$[Rx,3t+1]=[PRx,3t]=[x,3t]$, and the factors
$e^{-2\pi i/3}$ and $\zeta$ cancel in the disk coordinate.
The formula is therefore well defined. A lift of $\tau$ gives
$t=\tau/3$ and $s=e^{2\pi i\tau/3}z$; changing the lift is exactly
the quotient relation. This proves that the formula and its inverse
are smooth, including at $z=0$. Its real Jacobian has positive
factor $3$ in the $t$ direction and a rotation in the disk.
The formula for $q$ follows by cubing $s$.

At $v=1$ in \eqref{eq:positive-native-boundary}, applying the inverse
deck generator gives $[Px,t,1]$. This is the endpoint identification
of $B_M$, and $qj=e^{2\pi iv}$. Under $\Phi^+$ the map is
$([x,3t+v],e^{-2\pi it})$. The two circle coordinates have matrix
$\left(\begin{smallmatrix}3&1\\-1&0\end{smallmatrix}\right)$
of determinant one. Solving these coordinates as before gives a
smooth inverse compatible with the mapping-torus relations.
The boundary orientation is therefore preserved. Setting the disk
coordinate equal to zero gives both maps in
\eqref{eq:positive-geometric-maps}.
\end{proof}

No reversal of the positive base circle was used. For comparison,
the map $\mathcal J:B_M\to B_{M^{-1}}$ given by
$\mathcal J([x,t,v])=[Px,t,1-v]$ satisfies
$j_+\mathcal J=j_-$ for the earlier quotient $N$. It reverses the
base circle. There cannot be a comparison between these two earlier
markings covering the identity on the oriented base and inducing
the identity on fibre homology: it would intertwine $P$ and $R$
by the identity. But $Pu_0=u_1$ and $Ru_0=u_2$.

\section{Integral maps for positive transport}

The cell model for the core of $N^+$ is
\begin{equation}\label{eq:Q-positive}
 Q^+_n=D_n\oplus D_{n-1},\qquad
 d_+(x,y)=((P_{n-1}-I)y,0).
\end{equation}
The fibre complex remains $C$ and its boundary mapping torus remains
$B$, with differential $(a,b)\mapsto((M-I)b,0)$.
Set $S=I+P+P^2=I+R+R^2$ in each exterior degree.

\begin{theorem}\label{thm:positive-chains}
With the interval-first cell orientation, the geometric maps
\eqref{eq:positive-geometric-maps} induce
\begin{equation}\label{eq:positive-FHK}
 \begin{split}
 F^+_n(x,y)&=(x,S_{n-1}y),\\
 H^+_n(x,y)&=(0,x),\\
 K^+_n(x,y,z,w)&=(x,S_{n-1}y+z).
 \end{split}
\end{equation}
They satisfy
\begin{equation}\label{eq:positive-identities}
 d_+F^+=0,\qquad d_+H^++H^+d_C=F^+(M-I),
 \qquad d_+K^+=K^+d_B.
\end{equation}
These are maps on the actual fibre, core and boundary cell complexes,
not merely maps on their homology.
\end{theorem}
\begin{proof}
The cover $f^+$ traverses three core intervals in the positive
direction. The successive transports are $I,P,P^2$, which give
$S$ on the suspended summand and the identity on fibre cells.
The homotopy $[x,3t+v]$ traverses one positive core interval on
a cell at $t=0$. Its chain is $\sigma x$, giving $H^+(x,0)=(0,x)$.
On the $t$-prism of a $k$-cell, subdivide the $(t,v)$ square along
integer values of $3t+v$. Each image lies in the product of a
permuted $k$-cell and a core interval, hence in the $(k+1)$-skeleton.
Its degree-$(k+2)$ cellular component is zero. This proves the second
formula on every summand. The subdivisions agree on faces, so these
are cellular representatives of the covering and its stated homotopy.
Combining their fibre and boundary-interval components gives $K^+$.

For the chain identities, $(P-I)S=S(P-I)=0$ and
$d_+H^+(x,y)=((P-I)x,0)=F^+(M-I)(x,y)$.
The equality for $K^+$ follows by applying these two identities to
the fibre and suspended boundary summands separately.
\end{proof}

All entries through degree four can be read without an implicit
exterior operation. The groups $D_0,D_3$ have rank one and
$P_0=P_3=1$; $D_1,D_2$ have rank three and
\[
 P_1=P_2=\begin{pmatrix}0&0&1\\1&0&0\\0&1&0\end{pmatrix},
 \qquad S_0=S_3=3,\qquad
 S_1=S_2=\begin{pmatrix}1&1&1\\1&1&1\\1&1&1\end{pmatrix}.
\]
Thus $Q^+$ has ranks $1,4,6,4,1$, with only $d_{+,2}$ and
$d_{+,3}$ nonzero, each using the displayed $P-I$ on its suspended
triple. In particular $F^+_4$ is multiplication by three, $H^+_3$
sends the top fibre cell $\omega$ to its core prism, and $H^+_4=0$.
This also specifies all maps in the requested adjacent degrees.

\begin{proposition}\label{prop:positive-splitting}
The positive boundary map has chain section
$s^+_n(x,b)=(x,0,b,0)$. Its kernel consists exactly of
$(0,y,-Sy,w)$, with differential $(y,w)\mapsto((P-I)w,0)$.
The map $(a,b)\mapsto(-a,b)$ identifies $Q^+[1]$ with this kernel.
Consequently $B\cong Q^+\oplus Q^+[1]$ as integral chain complexes.
\end{proposition}
\begin{proof}
The formula for $K^+$ gives $K^+s^+=I$. Since the boundary of
its third component is $(P-I)b$, $s^+$ commutes with boundaries.
The equation $K^+=0$ forces $x=0,z=-Sy$. Its boundary has
second component $(P-I)w$ and third component zero, because
$(P-I)S=0$. This proves the kernel formula. On $Q^+[1]$ the
differential is $-d_+$; multiplication by $-1$ on its first
summand changes this to the displayed kernel differential.
All changes of basis are integral and invertible.
\end{proof}

The section is the actual normal-circle section $z=1$ in
$G_P\times S^1_z$: a positive core interval traverses the positive
boundary interval, with no inverse transport. The kernel carries
the extra normal circle. The image of $P-I$, like that of $R-I$,
is the full sum-zero lattice in each cyclic triple, and its kernel
is generated by the sum. Hence the proof of
\eqref{eq:Q-homology} applies with $P$ in place of $R$: explicitly
$H_n(Q^+)=\coker(P_n-I)\oplus\ker(P_{n-1}-I)$, with the same
integral generators $1,[u_0],[w_0],\omega$ and $1,U,W,\omega$.
The covering uses exactly the same $S$. Every homology matrix in
\eqref{eq:F-homology} therefore gives $F^+_*$ in these stated bases.
Its kernels are the indicated sum-zero lattices, and its only
nonzero cokernels are $\mathbb Z/3\mathbb Z$ in degrees one and four.
The groups of $N^+$ are torsion-free, of ranks $1,2,2,2,1$.
The boundary groups are torsion-free, of ranks $1,3,4,4,3,1$,
and $K^+_*$ is the split projection. Dualizing these free groups
gives specialization indices $3,1,1,3$ and degree-one image
$\mathbb Ze_3^*\oplus3\mathbb Ze_4^*$ in the unchanged fibre marking.

\section{The affine comparison required by a prescribed filling}

A marking of homology does not specify an origin of a torus.
In period coordinates an affine transformation has the form
$x\mapsto Ax+b$, where $b$ is a point of the torus. Changing the
origin by $c$ changes its translation to $b+(A-I)c$.
It does not change its homology matrix. A comparison of fillings
with a chosen boundary section must retain that section as well
as the integral periods.

\begin{proposition}\label{prop:marked-comparison}
Let a proper smooth torus group family over the $s$ disk have
the marked periods \eqref{eq:period-lattice}. Suppose its original
deck transformation over $s\mapsto\zeta s$ preserves its zero
section, has matrix $A=T^{-1}$, and the chosen origin on the
punctured quotient is the descent of that zero section. Twist the
deck transformation by the section $e_4(s)/3$. Then its quotient
is smoothly identified with $N^+$, with the same positive fibre
marking and boundary map \eqref{eq:positive-native-boundary}.
In particular \eqref{eq:Q-positive} and
\eqref{eq:positive-FHK} give its integral filling data.
\end{proposition}
\begin{proof}
The periods trivialize the real Lie algebras and exponentiate to
smooth group identifications with $\mathbb R^4/\mathbb Z^4$.
The deck matrix is locally constant, hence $A$ throughout the disk.
Preservation of zero makes its affine translation zero. The twisted
action is therefore $(x,s)\mapsto(Ax+e_4/3,\zeta s)$.
In the oriented $u$ basis this is \eqref{eq:positive-quotient}.
On the lifted positive boundary path $s=e^{2\pi iv/3}$,
the logarithmic translation
$(\log s)/(2\pi i)\,e_4$ is $ve_4/3$. It vanishes at $v=0$
and gives exactly \eqref{eq:positive-native-boundary}.
All subsequent geometric and chain comparisons have already been
constructed for this quotient. No reversal or change of periods is
involved.
\end{proof}

The zero-section hypothesis is substantive. To see its effect while
retaining the complex torus \eqref{eq:period-lattice}, let
$b\in\{0,1,2\}$ and define three holomorphic deck actions
\begin{equation}\label{eq:affine-ambiguity}
 \rho_b(x,s)=(Ax+be_4/3,\zeta s).
\end{equation}
All have order three, all have the same linear action on the full
integral lattice, and all induce the same action on the elliptic
coordinate. The descended punctured families are isomorphic as
marked torus families. In fact, on the universal logarithm cover put
\[
 \sigma_b(s)=\frac{b\log s}{2\pi i}\,e_4\in\mathbb C^2/\Lambda.
\]
A change of logarithm adds the integral period $be_4$, so this is a
single-valued holomorphic torus section on the punctured disk.
Since $Ae_4=e_4$, it satisfies
$\sigma_b(\zeta s)=A\sigma_b(s)+be_4/3$.
Translation by $\sigma_b$ consequently intertwines $\rho_0$ and
$\rho_b$, induces the identity on homology, and at $s=1$ fixes the
marked fibre point. It carries the descended zero section to a
section of the other punctured quotient.

Nevertheless twisting each action by the same specified $a=e_4/3$
produces translation $(b+1)e_4/3$. The cases $b=0,1$ are free,
whereas $b=2$ has fixed points on the central fibre: the translation
is then an integral period and the origin is fixed by $A$.
Thus positive integral monodromy, the elliptic deck calculation,
the increment $a$, and even an identified punctured family do not
imply freeness or select an affine filling. The examples establish
this precise failure of inference; they do not identify any one
of these alternative descents with a separately prescribed geometric
construction.

There is a further ambiguity even when the affine central action is
fixed. For an integer $r$, the formula
\[
 L_r([x,t,v])=[x,t+rv,v]
\]
defines a diffeomorphism of $B_M$ covering the identity on the
positive base and restricting to the identity on its marked fibre
at $v=0$. The endpoint condition holds because $r$ is integral
and $M$ fixes the $t$ circle. However the two boundary maps $j$
and $jL_r$ have distinct filling effects. Their retractions are
$[x,3t+v]$ and $[x,3t+(3r+1)v]$.
On the loop with $x=t=0$, they give respectively $[\tau]$ and
$(3r+1)[\tau]$ in $H_1(G_P;\mathbb Z)$, where $[\tau]$ has
infinite order. Thus an unspecified change of section along the
meridian changes a genuine attachment map while preserving the
linear monodromy and the marked starting fibre.

